\documentclass[preprint]{aastex631}
\usepackage{graphicx}
%\usepackage[authoryear,round,longnamesfirst]{natbib}
\bibliographystyle{aasjournal}
%\usepackage{natbib}
%\bibliographystyle{natbib}
%\bibliographystyle{unsrtnat}
%\bibliographystyle{apj}

\usepackage{mathtools}
%\usepackage{adjustbox}
\usepackage[mathlines]{lineno}
\linenumbers
\usepackage{xspace}
\graphicspath{ {./Figures/} }

\setlength\parindent{20pt}

\newcommand{\parti}[2]{\frac{\partial {#1}}{\partial {#2}}}

\usepackage{color}
\newcommand{\xxx}[1]{{\color{red}{#1}}}
\newcommand{\todo}[1]{{\color{blue}{\vspace{0.25cm}\noindent TO DO:#1}}}
\newcommand{\link}[1]{{\small{\texttt{\underline{#1}}}}}
\def\etal{{\it et al.\thinspace}}
\def\eg{{\it e.g., }}
\def\etc{{\it etc.\xspace}}
\def\ie{{\it i.e., }}
\def\cf{{\it c.f., }}
\def\gsim{~\rlap{$>$}{\lower 1.0ex\hbox{$\sim$}}}
\def\lsim{~\rlap{$<$}{\lower 1.0ex\hbox{$\sim$}}}
\def\wpsqm{{W/m$^{-2}$}}

\def\acen{{$\alpha$~Cen}}
\def\tess{{\it TESS}\xspace}
\def\kepler{{\it Kepler}\xspace}
\def\jwst{{\it JWST}\xspace}
\def\spitzer{{\it Spitzer}\xspace}
\def\aqua{{\it AQUA}\xspace}
\def\topex{{\it TOPEX/POSEIDON}\xspace}
\def\jason{{\it Jason}\xspace}

\def\msun{{M_\odot}}
\def\rsun{{R_\odot}}
\def\lsun{{L_\odot}}
\def\mearth{{M_\oplus}}
\def\rearth{{R_\oplus}}

\def\vplanet{\texttt{\footnotesize{VPLanet}}\xspace}
\def\atmesc{\texttt{\footnotesize{AtmEsc}}\xspace}
\def\binary{\texttt{\footnotesize{BINARY}}\xspace}
\def\distorb{\texttt{\footnotesize{DistOrb}}\xspace}
\def\distres{\texttt{\footnotesize{DistRes}}\xspace}
\def\distrot{\texttt{\footnotesize{DistRot}}\xspace}
\def\eqtide{\texttt{\footnotesize{EqTide}}\xspace}
\def\flare{\texttt{\footnotesize{FLARE}}\xspace}
\def\galhabit{\texttt{\footnotesize{GalHabit}}\xspace}
\def\galchem{\texttt{\footnotesize{GalChem}}\xspace}
\def\magmoc{\texttt{\footnotesize{MagmOc}}\xspace}
\def\poise{\texttt{\footnotesize{POISE}}\xspace}
\def\radheat{\texttt{\footnotesize{RadHeat}}\xspace}
\def\spinbody{\texttt{\footnotesize{SpiNBody}}\xspace}
\def\stellar{\texttt{\footnotesize{STELLAR}}\xspace}
\def\thermint{\texttt{\footnotesize{ThermInt}}\xspace}
\def\sporbres{\texttt{\footnotesize{SpOrbRes}}\xspace}
\def\vsp{\texttt{\footnotesize{VSPACE}}\xspace}
\def\mp{\texttt{\footnotesize{multi-planet}}\xspace}
\def\alabi{\texttt{\footnotesize{alabi}}\xspace}

\def\nbody6{\texttt{\footnotesize{NBODY6++}}\xspace}
\def\mercury{\texttt{\footnotesize{MERCURY}}\xspace}
\def\approxposterior{\texttt{\footnotesize{APPROXPOSTERIOR}}\xspace}
\def\emcee{\texttt{emcee}\xspace}
\def\dynesty{\texttt{dynesty}\xspace}
\def\otis{\texttt{\footnotesize{OTIS}}\xspace}
\def\rocke{\texttt{\footnotesize{ROCKE3D}}\xspace}
\def\deeptides{\texttt{\footnotesize{DEEPTIDES}}\xspace}

\def\k40{{$^{40}$K}}
\def\th232{{$^{232}$Th}}
% \def\uran238{{$^{238}$U}} Hopeless! Don't try to remake these!
% \def\uran235{{$^{235}$U}}
\def\coo{{CO$_2$}\xspace}
\def\water{{H$_2$O}}

\def\apjs{{\it Astrophys.~J.~Supp.}}
\def\apj{{\it Astrophys.~J.}}
\def\psj{{\it Plan.~Sci.~J.}}
\def\apjl{{\it Astrophys.~J.~Lett.}}
\def\aj{{\it Astron.~J.}}
\def\aap{{\it Astron.~\&~Astrophys.}}
\def\nat{{\it Nature}}
\def\mnras{{\it Mon.~Not.~R.~Astron.~Soc.}}
\def\icarus{{\it Icarus}}
\def\pasp{{\it Proc.~Astron.~Soc.~Pac.}}
\def\grl{{\it Geophys.~Res.~Lett.}}
\def\gca{{\it Geochem.~Cosmochem.~Acta}}
\def\jgr{{\it J.~Geophys.~Res.}}
\def\cemda{{\it Cel.~Mech.~Dyn.~Astron.}}
\def\asbio{{\it AsBio.}}
\def\scirep{{\it Sci.~Rep.}}

%\usepackage{hyperref}
 
\begin{document}

\title{The XUV Evolution of GJ 1132}

\author[0000-0001-6487-5445]{Rory Barnes}
\affiliation{Department of Astronomy, University of Washington, 3910 15th Avenue NE, Seattle, WA 98195, USA} 
\email{rory@astro.washington.edu}

\author[0000-0002-8341-0376]{Laura N. R. do Amaral}
\affiliation{School of Earth and Space Exploration, Arizona State University, 781 Terrace Mall, Tempe, AZ 85287, USA}

\author[0000-0002-0637-835X]{James R. A. Davenport}
\affiliation{Department of Astronomy, University of Washington, 3910 15th Avenue NE, Seattle, WA 98195, USA} 

\author[0000-0002-0637-835X]{Scott Engle}
\affiliation{Villanova University, Dept. of Astrophysics and Planetary Science, 800 E. Lancaster Ave., Villanova, PA 19085, USA}

\author[0000-0002-7961-6881]{Jessica Birky}
\affiliation{Department of Astronomy, University of Washington, 3910 15th Avenue NE, Seattle, WA 98195, USA} 

\author[0000-0002-0637-835X]{Evgenya Shkolnik}
\affiliation{School of Earth and Space Exploration, Arizona State University, 781 Terrace Mall, Tempe, AZ 85287, USA}

\author{Megan Gialluca}
\affiliation{Department of Astronomy, University of Washington, 3910 15th Avenue NE, Seattle, WA 98195, USA} 

\begin{abstract}
    We investigate the plausible range of XUV luminosity evolution of the planet-hosting star GJ 1132 to infer the total incident XUV on the short-period Earth-sized transiting planet GJ 1132 b. We included all observational uncertainties, compared two quiescent XUV luminosity models, and modeled the XUV luminosity from flares based on \tess data and a re-analysis of \kepler stars. We find that GJ 1132 has a low rate of flaring today, which suggests the M4 star is many Gyr old, a hypothesis bolstered by our inferred ages from both gyrochronology and quiescent XUV modeling. We then inferred the distribution of the cumulative XUV flux that planet b has intercepted. All tested models peak at a value that is 400-500 times larger than that of modern Earth, confirming that this planet is a good candidate to search for permanent atmospheric loss. We also find that an empirical XUV model for M dwarfs predicts a distribution that is $\sim$3 times narrower than the extrapolation of a solar twin XUV model. Flares contribute about 10\% of the total XUV luminosity, which, while modest, is enough to place planet b firmly on the  atmosphere-free side of the so-called ``cosmic shoreline''. Should GJ 1132 be typical of M dwarf planet hosts, then the uncertainty in short-period planets' cumulative XUV fluxes are about a factor of 3.  
\end{abstract}

\section{Introduction\label{sec:intro}}

\jwst is now searching for atmospheres of terrestrial exoplanets on short-period orbits, but it remains unclear how persistent these atmospheres may be \citep[see, \eg][]{Kreidberg19,Crossfield22,Xue2024,Mansfield24,Wachiraphan25}. In such tight quarters, the high energy radiation from the star can remove atmospheres by accelerating atoms to escape velocity, depleting the atmosphere over time. Thus, characterizing the historical incident XUV (1-1000\AA) flux of planet-hosting stars can provide critical insight into the longevity of short-period terrestrial planet atmospheres \citep{ZahnleCatling17}. Here we examine the XUV history of the planet host GJ 1132, an apparently bright star with at least two planets, one of which, planet b, is an optimal target for \jwst observations \citep{BertaThompson15}. Our ultimate goal is to predict the total number of XUV photons that have struck the planet over its lifetime.

While obtaining current X-ray luminosities, which can be scaled to XUV luminosities, is relatively straight-forward \citep{MoralesCalderon2024}, the current XUV luminosity provides minimal constraints on the historical XUV luminosity. Instead, models for the XUV evolution are assumed based on observations of stars with ages determined by association, \eg stars in open clusters or with white dwarf companions \citep{Engle24}. While imperfect, these methods currently offer the best path to estimating a star's XUV history and hence the cumulative XUV flux on a planet.

In general the XUV luminosity may be divided into quiescent and flaring components, which do not necessarly evolve in the same way. Several studies have developed empirical models for the quiescent  \citep{Ribas05, Engle24,France25} and flaring XUV evolution over time \citep{Loyd18,Davenport19,Ilin19,Guenther20,Feinstein20}. We will exploit these models to constrain GJ 1132's quiescent and flaring XUV evolution based on recent observation of GJ 1132's X-Ray luminosity and \tess data.

While numerous studies have computed a reasonable estimate of the cumulative XUV flux that these planets have intercepted, few, if any, have provided uncertainties. Given the relatively large uncertainties in age of most field stars, as well as uncertainty in parameters of XUV evolutionary models, the currently published estimates may be misleading. To begin to rectify the situation, here we introduce a framework to generate Bayesian posteriors for the cumulative XUV flux received by GJ 1132 b over its lifetime. Rather than attack the population of stars hosting short-period planets, we begin with this high priority target to identify the best methodology to infer uncertainties in the cumulative XUV flux on a planet. As shown below, the approach can easily be applied to other stars, assuming the observational data exist, to provide statistical constraints on any planet's cumulative XUV flux.

Table \ref{tab:starprops} gives parameters for the host star, with values taken from \citep[][hereafter Xue24]{Xue2024}, \citep[][hereafter MC24]{MoralesCalderon2024}, and \citep[][hereafter Bonfils18]{Bonfils2018}. These studies reveal an M4 dwarf star that appears to be relatively quiet today. While a calm chromosphere bodes well for robust detections of atmospheric or surface features with \jwst, it also makes it difficult to constrain the current flaring rate and age. Nonetheless, as we show below, the set of observations can be combined with modern statistical methods to infer the total XUV luminosity over time and hence the cumulative XUV flux on planet b.

\begin{table}
\centering
\caption{Literature Values for Selected GJ 1132 Properties}
\begin{tabular}{lcl}
\hline\hline
Property & Value & Ref.\\
\hline
Mass & $0.1945^{+0.0048}_{-0.0046}~\msun$ & Xue24\\
Total Luminosity & $0.00477^{+0.00036}_{-0.00026}~\lsun$ & Xue24\\
Radius & $0.2211^{+0.0069}_{-0.0081}~\rsun$ & Xue24\\
Density & $25.3^{+3.1}_{-2.2}$ g/cm$^3$ & Xue24\\
Effective Temperature & $3229^{+78}_{-62}$ K & Xue24\\
X-Ray Luminosity & $(9.96\pm2.95) \times 10^{25}$ ergs/s & MC24\\
Rotation Period & $122.3^{+6.0}_{-5.0}$ days & Bonfils18\\
\hline
\end{tabular}
\label{tab:starprops}
\end{table}

In the next section, we present the results of our search for flares on GJ 1132 in \tess data. Then we present a reanalysis of the flare data set from \citep{Davenport19} that takes into account non-detections.  In $\S$~\ref{sec:models} we present summaries of the physical models we used to calculate XUV evolution of GJ 1132. In $\S$~\ref{sec:stats}, we describe how we estimate uncertainties in the different aspects of the experiment. In $\S$~\ref{sec:results}, we present the results of the numerical modeling and estimate the cumulative XUV flux on GJ 1132 b for different assumptions. Finally, in $\S$~\ref{sec:conclusions}, we summarize our findings.


\section{Physical Models\label{sec:models}}

We use the planetary system evolution code \vplanet~\citep{Barnes20} to simulate the evolution of the star for the age of the system. This code includes quiescent and flaring models that predict the bolometric and XUV luminosity of the star, $L_{bol}$ and $L_{XUV}$, respectively. In the next three subsections we present the specific models we use in this study.

%The scripts to run the simulations and generate the figures is publicly available.\footnote{\href{https://github.com/VirtualPlanetaryLaboratory/vplanet/tree/master/examples}{\link{https://github.com/RoryBarnes/GJ1132}}}

\subsection{Structure\label{sec:starstructure}}

We use the \cite{Baraffe15} stellar evolution models to calculate $L_{bol}$, effective temperature $T_{eff}$, and radius $R_*$ as a function of time. This model depends only on the stellar mass and metallicity. As GJ 1132's metallicity is close to solar at [Fe/H] = -0.12, we use their solar metallicity track to model this star.

\subsection{Quiescent XUV Luminosity\label{sec:quiescencemodel}}

For the quiescent evolution, we consider two models: the ``solar twin'' model \citep[][hereafter R05]{Ribas05} and the ``empirical M dwarf'' (EMD) model \citep[][]{EngleGuinan23,Engle24}. We describe these two models in the following subsections.

\subsubsection{The Solar Twin Model\label{sec:solartwin}} 

The solar twin  is derived from 5 stars with masses and compositions that are nearly identical to the Sun, and so are significantly different from GJ 1132. Nonetheless, this model has been widely applied to M dwarfs \citep[\eg][]{ZahnleCatling17,Mordasini20,EveOwen25,Fleming20}, so we use it here as well. In particular, we use the parameterization described in \cite{LugerBarnes15}, which consists of 5 variables: age $t$, the ``saturation time'' $t_{sat}$, the ``saturation fraction'' $f_{sat}$, the post-saturation power law exponent $\beta_{XUV}$, and the mass $m$ (which sets the bolometric luminosity): 

\begin{align}
\label{eq:lxuv}
\frac{L_{XUV}^{quiescence}}{L_{bol}} = \left\{
        \begin{array}{lcr}
            f_\mathrm{sat} &\ & t \leq t_\mathrm{sat} \\
            f_\mathrm{sat}\left(\frac{t}{t_\mathrm{sat}}\right)^{-\beta_\mathrm{XUV}} &\ & t > t_\mathrm{sat}.
        \end{array}
    \right.
\end{align}



\subsubsection{Empirical M Dwarf Model\label{sec:empirical}} 

The EMD model consists of two-part linear relationship that is parameterized by the logarithm of the age in Gyr ($\equiv\tau$). The the age-rotation period ($P_{rot}$) is:
\begin{align}
\tau = \left\{
            \begin{array}{lcr}
                a_{rot}P_{rot} + b_{rot} &\ & P_{rot} < d_{rot} \\
                a_{rot}P_{rot} + b_{rot} + c_{rot}(P_{rot} - d_{rot}) &\ & P_{rot} \geq d_{rot},
        \end{array}
    \right.
    \label{eq:englerot}
\end{align}
when $P_{rot}$ is expressed in days. The XUV-age relationship is described by the following equation:
\begin{align}
\frac{L_{XUV}^{quiescence}}{L_{bol}} = \left\{
            \begin{array}{lcr}
                a_{xuv}\tau + b_{xuv} &\ & \tau < d_{xuv} \\
                a_{xuv}\tau + b_{xuv} + c_{xuv}(\tau - d_{xuv}) &\ & \tau \geq d_{xuv},
        \end{array}
    \right.
    \label{eq:englexuv}
\end{align}
The coefficients are provided in Table~\ref{tab:emd}. Note that for the XUV luminosity, the fits are only valid for M2 -- M6.5 stars and assumes a bandpass of 5--1700\AA, whereas the rotation period fit is valid for M4 -- M6.5 stars. Below we ignore differences between the XUV bandpass widths.

% To obtain the cumulative XUV flux that impinged upon a given planet for this model, we first determine the distribution of ages consistent with the rotation period, Eq.~(\ref{eq:englerot}) and then integrate Eq.~(\ref{eq:englexuv}) for each plausible age. In other words,
% \begin{equation}
%     F_{XUV}^{cumulative} = \frac{1}{4\pi a^2}\Bigg(\int_0^{d_{xuv}}(a_{xuv}\tau' + b_{xuv})d\tau' + \int_{d_{xuv}}^{\tau}(a_{xuv}\tau' + b_{xuv} + c_{xuv}[\tau' - d_{xuv}])d\tau'\Bigg),
% \end{equation}
% assuming a circular orbit.

\begin{deluxetable*}{c|cc|cc}
\tablecaption{Coefficients for the EMD Model \label{tab:emd}}
\tablehead{\colhead{}  & \multicolumn{2}{c}{XUV}
           & \multicolumn{2}{c}{Rotation}\\
            \hline
           \colhead{Parameter} & \colhead{Best Fit} & \colhead{Uncertainty}  
           & \colhead{Best Fit} & \colhead{Uncertainty}}
\startdata
\hline
a & -0.1456 & 0.0911 & 0.0251 & 0.0018\\
b & -2.8876 & 0.0439 & -0.1615 & 0.0303\\
c  & -1.8187 & 0.2412 & -0.0212 & 0.0018\\
d & 0.3545  & 0.0604 & 25.4500 & 1.9079\\
\hline
\enddata
\end{deluxetable*}

\subsection{XUV from Flares\label{sec:flaremodel}}

Our approach to modeling the XUV output from stellar flares follows exactly the model described in \citet{Amaral22}. Stellar flares are typically described in terms of the flare frequency distribution (FFD), a power law that describes the frequency of flares per unit energy $\nu$: 
\begin{equation}
\label{eq:ffd}
    \log(\nu) = a\log(E) + b,
\end{equation}
where $E$ is the flare energy, $a$ is the power law index (slope in log-log visualizations), and $b$ is a scaling factor (the $y$-intercept in log-log visualizations). The slope and $y$-intercept of the FFD is observed to change with time \citep[\eg][]{Davenport19}, so in general the cumulative XUV from flares must be modeled according to the time rates of change of $a$ and $b$.

A fraction of the energy from each flare is emitted in the XUV and we use the conversion factors from
\citet[][Table 2]{OstenWolk15}. Then the instantaneous XUV luminosity from flares is given by
\begin{equation}
       L_{XUV}^{flares} = \int_{E_{min}}^{E_{max}} \nu(E_{XUV})\ dE,
\label{eq:LXUVFlare}
\end{equation}
where $E_{min}$ and $E_{max}$ are the minimum and maximum flare energies considered, and $E_{XUV}$ is the portion of the flare energy emitted in the XUV. Thus, with a model for the FFD with time, we can compute the XUV luminosity from flares. The total XUV luminosity is then simply $L_{XUV} = L_{XUV}^{quiescence} + L_{XUV}^{flares}$. For more details on the approach, consult \citet{Amaral22}.


% \begin{table}
% \centering
% \caption{\cite{Davenport19} Model Parameters and Uncertainties}
% \begin{tabular}{lcc|lcc}
% \hline\hline
% Property & Mean & Uncertainty & Property & Mean & Uncertainty\\
% \hline
% $a_1$ & -0.07 & 0.11 & $b_1$ & 2.01 & 3.84\\
% $a_2$ & 0.79 & 0.40 & $b_2$ & -25.15 & 13.58\\
% $a_3$ & -1.06 & 0.36 & $b_3$ & 33.99 & 12.35\\
% \hline
% \end{tabular}
% \label{tab:dav19params}
% \end{table}



\section{Statistical Methods\label{sec:stats}}

To constrain the outcomes of our models we use a mix of Bayesian inference for model parameters that are directly constrained by observational data, and ``convergence modeling'' that convolves Bayesian posteriors with reasonable priors to derive distributions of all model parameters. This mix takes advantage of \alabi's machine learning capabilities as well as simplifies the problem for the unknown parameters. The disadvantage of this approach is that it relies on priors that are not rigorously constrained, in other words the shapes of the priors for many variables that are critical for total devolatization are chosen based on the intuition of the experimenter. Whenever possible, we point out areas where these intuitively informed priors may impact final probability distributions. We follow the approach laid out in \citep{Gialluca24}.

\subsection{Bayesian Inference\label{sec:alabi}}

We use the code \alabi~\citep{Birky25} to infer posteriors of the quiescent stellar evolution. This code trains only on the likelihood, so is unable to constrain parameters that do not impact the likelihood. We do not use the planetary data in conjunction with the stellar data to infer model parameters. We find that in many cases the planetary constraints will strongly influence the stellar posteriors in ways that seem very unlikely. Moreover, we do not vary the planetary mass or radius over the simulation, so in this model, the planet does not perturb the star.

We infer the values of the model parameters via the procedure described in \cite{Fleming20} in which a machine learning model derives their posteriors. This methodology is described in more detail in $\S$~\ref{sec:alabi} and \cite{Birky25}. To obtain the cumulative XUV flux received by planet b, we use the methodology described in $\S$~\ref{sec:vconverge} and \cite{Gialluca24}.

\subsection{Convergence Modeling\label{sec:vconverge}}

To calculate dependent parameter distributions we performed sets of \vplanet simulations with initial conditions ranges set to posteriors determined via \alabi when possible or with well-defined priors when not. As described in \cite{Gialluca24}, after an initial set of trials is performed, batches of additional trials are performed and the distributions of specific final variables are compared between iterations until convergence. For example, an initial run may consist of 1000 trials, then batches of 500 trials are run until the distributions of the amount of water, \coo, and oxygen in the atmosphere are changing by less than 1\% over three consecutive batches. We use the K-L divergence test to compare the distributions, with convergence requiring the K-L statistic to remain below 0.004 for all parameters three times in a row.

This two-tiered approach saves significant computational time compared to standard MCMC methods. Bayesian inference requires $<1$\% of the resources as \texttt{emcee} and \texttt{dynesty}. It should be noted, however, that \alabi does depend on an MCMC solver such as those two. So with the directly constrainable parameters efficiently obtained, the most efficient way to compute the probability distribution for all parameters is to then run an additional set of runs to convergence. This step requires more computer resources than \alabi, but the total investment is still only of order 1\% of an MCMC run with \vplanet.

\section{Reanalysis of Kepler Flares\label{sec:kepler}}

After testing multiple flare evolution models \citep{Loyd18,Davenport19,Ilin19,Guenther20,Feinstein20}, we ultimately found the \citep{Davenport19} matched the observations best and applied it to GJ 1132. The Davenport et al.~model is a six-parameter fit in which the flare frequency follows a power law in terms of flare energy, and the slope and $y$-intercept evolve with time $t$. The model relates the flare frequency $\nu$ to the slope and $y$-intercept as
\begin{equation}
\label{eq:ffd_a}
    a = a_1\log t + a_2 m + a_3,
\end{equation}
and
\begin{equation}
\label{eq:ffd_b}
    b = b_1\log t + b_2 m + b_3,
\end{equation}
where $m$ is the stellar mass. For more information on the details of the model and its motivation, we refer the reader to \citep{Davenport19}.

As the original Davenport et al.~model did not include data points in which flares when fitting for $a_1, ..., b_3$, it overpredicted the FFD compared. We thus performed a new MCMC analysis with \emcee that included the non-detections in order to generate a more accurate model for GJ 1132. We used the same data set as Davenport et al., which consists of X stars and nearly 500,000 data points. The resulting corner plot is shown in Fig.~\ref{fig:ffd_corner}, which is similar to the results in \citet{Davenport19}, including some tight degeneracies in the covariances. The resulting best fits and 1$\sigma$ uncertainties are shown in Table \ref{tab:ffd}. Note that this is the global fit, so the parameters apply to all stars, not just GJ 1132.

\begin{figure}
\centering
    \includegraphics[width=\textwidth]{Figures/CornerVariableSlope.png}
    \caption{Corner plot for the six-parameter fit for the \citet{Davenport19} FFD model. The blue contours denote the 16\%, 50\%, and 84\% confidence intervals, with additional, lower-likelihood points shown by dots. The marginal posteriors along the diagonal include vertical dashed lines that denote the median and 1$\sigma$ uncertainties in the model parameter distributions.}
    \label{fig:ffd_corner}
\end{figure}

\begin{table}
\centering
\caption{Fit Parameters for FFD Evolution from \kepler Data}
\begin{tabular}{ll|ll}
\hline\hline
$a_1$ & $-0.148 \pm 0.000529$ & $b_1$ &  $4.69 \pm 0.0180$\\
$a_2$ & $0.517 \pm 0.00244$ & $b_2$ & $-16.45 \pm 0.0824$\\
$a_3$ & $-0.618 \pm 0.00231$ & $b_3$ & $19.446 \pm 0.0779$\\
\hline
\end{tabular}
\label{tab:ffd}
\end{table}

Figure \ref{fig:ffd_compare} compares the original fit to the new one for a 0.5 $\msun$ star, which is a mass that Davenport et al. focused on. We find that the predicted FFDs are similar, but the new fits predict lower overall flare rates as expected. At younger ages, the slopes are flatter, indicating higher energy flares are more common. We stress that, as in Davenport et al., the fit does not match all stars well, but does predict FFDs for many stars. We also tried a three-parameter fit in which $b = -1$ for all cases, but found that it did not perform significantly better. We therefore adopt the posteriors from Fig.~\ref{fig:ffd_corner} for our modeling of GJ 1132. Although there are degeneracies in the model, our procedure will not be affected by them, see $\S$~\ref{sec:stats}.

\begin{figure}
\centering
    \includegraphics[width=\textwidth]{Figures/ffd_comp.png}
    \caption{Comparison of the new FFD model (solid line) to the original \citet{Davenport19} FFD model (dashed lines) for 4 different ages.}
    \label{fig:ffd_compare}
\end{figure}

Finally, we note that the current flare rate can also provide a constraint on age. Once the slope, $y$-intercept, and stellar mass are measured, an age probability distribution can be inferred. Thus, the FFD provides a third mechanism to constrain the system's age.

\section{Current Flare Rate\label{sec:obsflare}}

FFDs can be difficult to measure for many stars, but the high cadence, long baseline observations of the \kepler and \tess spacecraft have enabled robust flare detections \citep{Davenport19,Guenther20}. Fortunately, GJ 1132 lies in a \tess field, and so we can use the \tess data to obtain an FFD for GJ 1132 today.

To derive the current FFD for GJ 1132, which can then be converted into an XUV luminosity, $L_{XUV}^{flares}$, we downloaded and analyzed all 5 sectors of available \tess data that include this star (123.31 days). In general, we find that the frequency of flares is quite low for this star, including no detections in Sectors 88 and 89 (Jan--Mar 2025). However, the lightcurves also show continuous low level variations ($\sim$0.1\%) that are above the noise and that may indicate constant low-level flaring. Figure \ref{fig:flareex} shows three example flares detected in \tess. 

We calculated the FFD by following the procedure described in \citet{Davenport19} and show the derived FFD in Figure \ref{fig:ffd}. The figure also shows the FFD for several other stars fro reference. The relatively low FFD for GJ 1132 indicates an age of at least 1 Gyr, but it is probably much older. We find $a = -0.32 \pm 0.77$ and $b = 8.54 \pm 24.34$. The large uncertainties reflect the small number of detected flares
and reveal that the current and historical FFD of GJ 1132 is very poorly constrained by observational data. 

\begin{figure}
\centering
    \includegraphics[width=\textwidth]{Figures/GJ1132_flares.pdf}
    \caption{Three flares of GJ 1132 observed by \tess. The blue points indicate times that are labeled as flaring and the used to compute the total energy of the flare.}
    \label{fig:flareex}
\end{figure}

\begin{figure}
\centering
    \includegraphics[width=\textwidth]{Figures/GJ1132_FFD_comp2.pdf}
    \caption{GJ 1132's FFD from \tess data (black points) as well as a linear regression fit to the data (black line). Additional FFDs are include for reference including the average of several clusters today (dashed red, dark blue, and orange lines) as well as the prediction of our flare model after an age of 8 Gyr (solid red, dark blue, and orange lines). The FFDs of two M dwarfs observed with \kepler, GJ1243 (age = ?) and GJ 4083 (age = ?), are shown with dotted pale blue lines.}
    \label{fig:ffd}
\end{figure}

\section{Current Quiescent XUV \label{sec:obsflare}Luminosity}

Our final step to calculate the cumulative XUV luminosity is to determine the current XUV luminosity, from which we can infer the values of the \citet{Ribas05} model parameters. As shown in Table~\ref{tab:starprops}, only the XUV luminosity has been directly measured, so we must rely on a scaling relationship to calculate the EUV luminosity, $L_{EUV}$. For this calculation, we use the \citet{SanzForcada25}:
\begin{equation}
    \log(L_{EUV}) = (0.821 \pm 0.041)C_X + (28.16 \pm 0.05),
    \label{eq:leuv}
\end{equation}
where
\begin{equation}
    C_X = \log(L_X) - 27.44
\end{equation}
when all luminosities are expressed in units of ergs/s. We first build the $L_{XUV}$ distribution via $10^4$ draws from the observed $L_X$ distribution, and then computing an $L_{EUV}$ value from it. The result is shown in the left panel of Fig.~\ref{fig:quiescentXUV}, where the best fit and uncertainty are $L_{XUV} = (2.74 \pm 0.83) \times 10^{-7} L_\odot$. Figure \ref{fig:quiescentXUV} shows that the uncertainties are well-modeled by a Gaussian.

While the absolute value of the XUV luminosity is valuable, the \citet{Ribas05} model requires the ratio of $L_{XUV}/L_{bol}$, see Eq.~\ref{eq:lxuv}. We thus created this distribution via $10^4$ independent draws from the normally distributed $L_{XUV}$ distribution and the asymmetric uncertainties in $L_{bol}$ and find $\log(L_{XUV}/L_{bol}) = -4.26 \pm 0.15$. The resulting distribution is shown in the right panel of Fig.~\ref{fig:quiescentXUV}, which shows some skew due to the asymmetric uncertainties in $L_{bol}$, especially at low values. Nonetheless, the uncertainties are close enough to a normal distribution that we adopt a symmetric Gaussian to represent the quiescent XUV luminosity.

\begin{figure}
    \includegraphics[width=0.47\textwidth]{Figures/lxuv_hist.pdf}
    \includegraphics[width=0.47\textwidth]{Figures/log_lxuv_lbol_hist.pdf}
    \caption{{\it Left:} The distribution of quiescent XUV luminosity of GJ 1132. {\it Right:} The distribution of the ratio of the XUV-to-bolometric luminosity of GJ 1132.}
    \label{fig:quiescentXUV}
\end{figure}


\section{XUV Evolution\label{sec:results}}

In this section we present the results of our stellar modeling, including XUV from quiescence and flares. We first present the evolution of the star, then cumulative XUV flux impinging on planet b.

\subsection{Stellar Evolution}

We begin by considering the \citet{Ribas05} model and use the $L_bol$ and $L_{XUV}/L_{bol}$ constraints to infer posteriors for its model parameters. Table \ref{tab:ribas} list the priors, expressed as a Gaussian, and bounds for the model parameters. The $t_{sat}$ prior was uninformative. Fig.~\ref{fig:XUVPosteriors} is the posteriors for the \citep{Ribas05} XUV model as generated by \texttt{alabi} and \texttt{VPLanet}. We used both the \texttt{emcee} and \texttt{dynesty} solvers and found very similar results, however the values for $t_{sat}$ are poorly constrained. 

\begin{table}
\centering
\caption{Priors and Bounds Used for Constraining the \citet{Ribas05} Model Parameters}
\begin{tabular}{l|cc|cc}
\hline\hline
Parameter & Mean & Standard Deviation & Lower Bound & Upper Bound\\
\hline
Mass [$M_\odot$] & 0.181 & 0.019 & 0.1 & 0.3\\
log($f_{sat}$ [Gyr] & -2.92 & 0.26 & -4.0 & -1.0\\
$t_{sat}$ [Gyr] & - & - & 0.1 & 5\\
$\beta_{XUV}$ & -1.18 & 0.31 & -2.0 & 0.0\\
\hline
\end{tabular}
\label{tab:ribas}
\end{table}

\begin{figure}
    \includegraphics[width=\textwidth]{Figures/gj1132posteriors.png}
    \caption{Corner plot for the \cite{Ribas05} model parameters. Blue curves are from \texttt{emcee}, red from \texttt{dynesty}, and grey are priors (see text).}
    \label{fig:XUVPosteriors}
\end{figure}


Figure \ref{fig:EngleAges} shows the distribution of ages predicted by the \cite{Engle24} model based on the observed constraint in the current rotation period. The distribution is well-described by a Gaussian distribution: $\tau = 0.76 \pm 0.22$, which corresponds to an age of $\sim$6 Gyr. Note that we assume the star cannot be older than 13 Gyr. Figure \ref{fig:example} shows 100 representative examples for each quiescent model plus flares. This figure clearly shows the narrower predictions of the \citet{Engle24} model compared to \citet{Ribas05}.

\begin{figure}
    \includegraphics[width=\textwidth]{Figures/EngleAgeHist.pdf}
    \caption{Distribution of GJ 1132 ages predicted by the \cite{Engle24} model, truncated at 13 Gyr. \xxx{Add Ribas predictions.}}
    \label{fig:EngleAges}
\end{figure}

\begin{figure}
    \includegraphics[width=\textwidth]{Figures/XUVEvol.pdf}
    \caption{Example XUV evolutions for GJ 1132. Both panels include flares, but the left panel assumes the quiescent model of \citet{Engle24} and the right assumes the model from \citet{Ribas05}.}
    \label{fig:example}
\end{figure}

% \begin{figure}
%     \includegraphics[width=\textwidth]{Figures/FFD_GJ1132.png}
%     \caption{Predicted evolution of GJ 1132's flare frequency distribution.}
%     \label{fig:ffdevol}
% \end{figure}

% Table \ref{tab:flareresults} shows the distribution of the \cite{Feinstein20} model parameters we obtained for the \tess data.

% \begin{table}
% \centering
% \caption{Flare Properties of GJ 1132}
% \begin{tabular}{lc}
% \hline\hline
% Property & Value\\
% \hline
% $f_1$ & $-0.31199 \pm 2 \times 10^{-4}$\\
% $f_2$ & -0.5 $\pm$ 0.7\\
% $f_3$ & 0 (fixed)\\
% $f_4$ & 30 $\pm$ 60\\
% \hline
% \end{tabular}
% \label{tab:flareresults}
% \end{table}

\subsection{Cumulative XUV Flux on GJ 1132 b}

With the distribution of XUV luminosity over time computed, we turn to its potential impact on planet b. Using the \cite{Ribas05} model with a saturation time of 100 Myr and a saturation fraction of $10^{-3}$, the cumulative quiescent flux the Earth has received is $F_{XUV,\oplus}^{Cumulative} = 9.8 \times 10^{15}$ W. We assume that Earth's distance from the relatively quiet Sun implies that the solar flare contribution to Earth's cumulative XUV flux is negligible.

%Figure \ref{fig:cumulativexuv} presents the probability distribution of planet b's $F_{XUV}^{Cumulative}$ based on the posteriors from Fig.~\ref{fig:XUVPosteriors}. 


Figure \ref{fig:shoreline} shows these XUV ranges in relation to the so-called ``cosmic shoreline,'' an empirical line inspired by energy-limited escape that separates bodies that possess significant atmospheres from those that don't \citep{ZahnleCatling17}. Among all our trials, only X\% of the Solar Twin model without flares lie on the atmosphere-expected side and that including flares move the lower limit well above the shoreline. Thus, it seems very likely that the Cosmic Shoreline model strongly favors an atmosphere-free world.

Figure \ref{fig:shoreline} and Table \ref{tab:results} show that the two quiescent XUV model predict mean values that differ by only $\sim$5\%. The empirical M dwarf model, however, produces a much narrower distribution of cumulative XUV fluxes than than the Solar Twin model. The 95\% confidence interval of the Ribas \etal model is 3.1 times larger than that of Engle. 

\begin{figure}
    \includegraphics[width=0.47\textwidth]{Figures/GJ1132b_CumulativeXUV_Multi.png}
    \includegraphics[width=0.47\textwidth]{Figures/CosmicShoreline.png}
    \caption{{\it Left:} The distribution of cumulative XUV fluxes planet b has received for different assumptions. {\it Right:} The distribution of XUV fluxes for GJ 1132 b in relation to the ``cosmic shoreline'' \citep{ZahnleCatling17}. The uncertainties shown are two sigma. The GJ 1132 b points are offset slightly in the $x$-axis to improve readability. The Solar System planets are shown for reference.}
    \label{fig:shoreline}
\end{figure}

\begin{table}
\centering
\caption{Cumulative XUV Flux on GJ 1132 b (Earth Units) }
\begin{tabular}{lcc}
\hline\hline
Model & Mean & 95\% Confidence Interval\\
\hline
Engle Only & 401.16 & [238.48, 564.09]\\
Engle w/Flares & 478.71 & [275.49, 662.63]\\
Ribas Only & 420.60 & [80.82, 1104.60]\\
Ribas w/Flares & 478.31 & [141.18, 1106.13]\\
\hline
\end{tabular}
\label{tab:results}
\end{table}


\section{Discussion\label{sec:discussion}}

We have performed a rigorous analysis of the history of XUV luminosity of GJ 1132, a nearby planet-hosting star whose innermost planetary companion is the target of ongoing JWST campaigns to detect and characterize a putative atmosphere. We are able to thoroughly constrain its quiescent and flaring XUV luminosity over time based on observations and find that planet b has likely intercepted over 100x more XUV photons than Earth. Combined with the known mass and radius, we find that the Cosmic Shoreline hypothesis \citep{ZahnleCatling17} predicts the planet will be devoid of an atmosphere, a prediction currently being checked via \jwst observations.

To the best of our knowledge, this work constitutes the first rigorous treatment of the uncertainties of the cumulative XUV flux for a terrestrial \jwst target. Our results reveal that the uncertainties in this value span a range of 3--10 depending on the model. Given that the EMD model is both calibrated to M dwarfs and produces a tighter probability distribution, we argue that it should be used in the future to constrain the XUV flux on short-period planets. With the first star characterized in this way, an obvious next step is to generate a catalog of XUV luminosities and the resultant planet fluxes to both compare to the ``cosmic shoreline'' and as inputs to probabilistic planetary volatile models \citep[\eg][]{KrissansenTottonFortney22,Gialluca24}

A critical challenge in modeling the cumulative XUV flux is the stellar age, which is notoriously difficult to measure for field stars. In this case, we obtained three different age estimates ($L_{XUV}/L_{bol}$, gyrochronology, FFD) that are all in agreement, mostly because the uncertainties of the individual measurements are so large. Nonetheless, that all methods are in agreement at least gives some hope that if only one method is available, it is reasonable to assume that the value is accurate, if not precise. Regardless, a uniform application of an age method is probably advisable when comparing planets.

Our comparison of the Ribas et al. and Engle XUV models has revealed that constraining the XUV history of field stars is no longer limited by model uncertainties. In Fig.~\ref{fig:ErrorCompare}, we plot the cumulative XUV flux for the two models with all uncertainties (black), model errors only (purple; stellar parameters held fixed at the best-fit and an age of 8 Gyr), and stellar errors only (red; model parameters held fixed at their best fit value). For the Ribas et al. model errors, we used the \dynesty posterior shown in Fig.~\ref{fig:XUVPosteriors}. Flaring is not included in these distributions. For the new Engle model, the distributions are very similar, suggesting the neither the stellar or model uncertainties are dominating the overall error budget. For the Ribas et al. model, however, the uncertainties in the model parameters clearly dominate. Thus, while the two models predict about the same mean cumulative XUV flux, the \citet{Engle24} model has clearly improved the precision of XUV modeling.

\begin{figure}
    \includegraphics[width=\textwidth]{Figures/GJ1132b_ErrorSourceComparison.png}
    \caption{Comparison of cumulative XUV flux distributions for the \citet{Engle24} model (left) and \citet{Ribas05} model (right). The red histograms only included uncertainties in the fundamental stellar parameters (mass and age); the purple histograms only included uncertainties in the model parameters; the black histogram includes both. The cosmic shoreline is shown in pale blue for reference.}
    \label{fig:ErrorCompare}
\end{figure}

The addition of flares to the total cumulative XUV flux represents a conceptual change to the cosmic shoreline model, so may not be appropriate to include. The original formulation is based solely on the quiescent XUV flux of the Sun, ignoring relatively infrequent flares that are likely fairly diluted by the time they reach the terrestrial planets of the Solar System. Although we find flares only provide about 10\% of the total XUV budget, should they be implicitly included in the cosmic shoreline model, including flares could be pessimistic in terms of atmosphere retention.

Our revised Kepler flare model predicts lower overall flare rates, see Fig.~\ref{fig:ffd}, but still predicts a significantly steeper and higher FFD for GJ 1132, see Fig.~\ref{fig:ffd_compare}. However, it is important to note that many \kepler flares possess the same morphology and the new fit includes large uncertainties, see Table~\ref{tab:ffd}.  Nonetheless, the inferred current FFD is still consistent with the \kepler analysis, as shown in Fig.~\ref{fig:FitCompare}, which shows that there are solutions to current FFD (pale blue points) that are very similar to the \kepler prediction (orange points). We thus conclude the results are consistent due to the large uncertainties in both fits.

\begin{figure}
    \includegraphics[width=\textwidth]{Figures/FitComparison.pdf}
    \caption{Comparison of the slope and $y$-intercept of the \kepler FFD fit at 8 Gyr to the range of values found from analyzing the GJ 1132's \tess lightcurve. The range of GJ 1132 values overlaps with the \kepler fit, demonstrating that GJ 1132's FFD is compatible with the \kepler data.}
    \label{fig:FitCompare}
\end{figure}

Although we have shown that GJ 1132 b sits firmly on the atmosphere-free side of the cosmic shoreline, that should not be taken to mean that it has no atmosphere, let alone that it has been permanently lost. While light elements such as H, He, and even O can be lost thermally \citep{Zahnle90}, the removal of carbon dioxide is very difficult \citep{KrissansenTottonFortney22}. CO$_2$ could still be lost mechanically through ablation by stellar energitic particles from coronal mass ejections (CMEs), which can strike a planet at relativistic speeds and transfer sufficient momentum to heavy molecules for them reach escape velocity. With the recent detection of the first CME beyond the Solar System \citep{Callingham25}, the community may be on the precipice of directly constraining the CME frequency distribution of solar-type stars. Once the stellar XUV and CME histories can be constrained, the rate of atmospheric loss of volatile can be properly calculated. 

The stellar energetic particles associated with CMEs drive non-thermal escape, which can be influenced by a magnetic field \citep{DriscollBercovici13,Gunell18}, but in principle the loss rate can be constrained. Since some planets will likely be magnetized, a full accounting of atmospheric retention must include the geodynamo, which is driven by convection in the outer envelopes of cores, which is set by the temperature difference between the core and mantle \citep{olson2006}. Thus, a geodynamo prediction must also include the mantle's thermal evolution, which is strongly influenced by the volatile evolution, both in terms of the radiative greenhouse warming of the atmosphere as well as the lubricating effect of water in the mantle that reduces the temperature needed for it to convect \citep{FoleyDriscoll16,McGovernSchubert89,Garcia25}. Thus, a realistic model of permanent atmospheric loss requires a core/mantle/crust/atmosphere/stellar model that couples all these pieces together. Future work should ultimately reach this stage because predictions for atmospheric loss clearly depend on the more stellar and planetary modeling than presented here. 

\section{Conclusions\label{sec:conclusions}}

We have analyzed the XUV luminosity and found that the planet has likely been hit by 400-500 times more XUV photons than Earth, putting its atmosphere at risk of total loss. We compared 2 quiescent XUV models and 1 flaring XUV model and find that all assumptions generate a cumulative XUV flux that the cosmic shoreline hypothesis predicts would result in an atmosphere-free world. We find that the \citet{Ribas05} and \citet{Engle24} models are consistent, but that the latter's uncertainties are about 3 times smaller, suggesting it is the preferred model for stars similar to GJ 1132.

Our framework can be applied to additional planets to help identify those planets that are likely to be devoid of an atmosphere or not. Crucially, by including uncertainties in the cumulative XUV flux, the interpretation of \jwst data in the context of the cosmic shoreline can be more rigorous. The research described herein motivates a comprehensive study of the uncertainties in the cumulative XUV flux on \jwst targets.

This research contributes to our understanding of how planets permanently lose their atmospheres, but does not yet produce a realistic prediction for the pathway to total volatile loss. To fully address the processes that lead to permanent atmospheric loss, stellar XUV models like those considered here must be stitched together with planetary models to create a new planetary system model that tracks planetary volatiles between mantle, crust, atmosphere, and space reservoirs, including the role of forced tidal heating in the planetary mantle. Ultimately, the synergy of theory and observation will mature to enable a full understanding of how planets devolatilize and the fate of terrestrial planets throughout the universe.


\vspace{1cm}

%\begin{acknowledgments}
\noindent{\bf Acknowledgments}\\
RB and JB acknowledge support from NASA grant No. 80NSSC23K0261.

\vspace{1cm}

\noindent{\it Software:} \texttt{alabi}, \texttt{astropy}, \texttt{bigplanet}, \texttt{dynesty}, \texttt{emcee}, \texttt{FFD}, \texttt{lightkurve}, \texttt{matplotlib}, \texttt{multiplanet}, \texttt{numpy}, \texttt{scipy}, \texttt{vconverge}, \texttt{vplanet},  \texttt{vspace}.

%that permits a prediction for the age, albeit with some unphysically large ages that must be manually excised from the distribution. 

\bibliography{gj1132,here}

\end{document}
